\subsection{Ideals and Generated Ideals}

Ideals are the subsets by which quotient rings can be formed, just as subspaces are the subsets by which quotient vector spaces can be formed.

\begin{definition} \label{def:ideal-and-generated-ideal} \hlt{Ideals and Generated Ideals}\\
Let $R$ be a ring.
\begin{enumerate}[label=(\roman*)]
\setlength{\itemsep}{0pt}
\item A nonempty subset $\mathcal{I}\subseteq R$ is a \hlt{two-sided ideal} if $a-b\in\mathcal{I}$ $\forall a,b\in\mathcal{I}$, and $ra,ar\in\mathcal{I}$ $\forall a\in\mathcal{I}$, $r\in R$.
\item If $R$ is commutative with identity, $S\subseteq R$, \hlt{ideal generated by $S$} is $\langle S\rangle\defeq\left\{\sum\limits_{j=1}^{n}r_js_j:n\geq1,\ r_j\in R,\ s_j\in S\right\}$.
\item \hlt{Finite generated ideal notation}: if $S=\{s_1,\ldots,s_n\}$, write $\langle s_1,\ldots,s_n\rangle\defeq\langle S\rangle$.
\end{enumerate}
\end{definition}

\begin{example} \label{ex:principal-ideal-polynomial-multiples} \hlt{Principal Ideals in Polynomial Rings}\\
In $\mathbb{F}[x]$, $\langle p(x)\rangle=\{q(x)p(x):q(x)\in\mathbb{F}[x]\}$ is the \hlt{principal} ideal generated by $p(x)$.
\end{example}

\begin{proposition} \label{prop:generated-ideal-smallest-containing-set} \hlt{Generated Ideals Are Smallest}\\
Let $R$ be a commutative ring with identity and let $S\subseteq R$.\\
Then $\langle S\rangle$ is an ideal, $S\subseteq\langle S\rangle$, and for every ideal $\mathcal{J}\subseteq R$, $S\subseteq\mathcal{J}\implies \langle S\rangle\subseteq\mathcal{J}$.
\begin{proof}
Let $x=\sum_{j=1}^{m}r_js_j$ and $y=\sum_{k=1}^{n}t_ku_k$ lie in $\langle S\rangle$.\\
Then $x-y=\sum\limits_{j=1}^{m}r_js_j+\sum\limits_{k=1}^{n}(-t_k)u_k\in\langle S\rangle$; $ax=\sum\limits_{j=1}^{m}(ar_j)s_j\in\langle S\rangle$ for all $a\in R$, so $\langle S\rangle$ is an ideal.\\
Since $s=1_Rs\in\langle S\rangle$, $S\subseteq\langle S\rangle$. If $\mathcal{J}$ is an ideal with $S\subseteq\mathcal{J}$, then $r_js_j\in\mathcal{J}$ and finite sums remain in $\mathcal{J}$, QED.\\
\hlt{Structural Commentary}: Generated ideals close a set under exactly the operations ideals must allow.\\
\hlt{Key Insight}: To force $S$ into an ideal, one must include all ring-linear combinations of elements of $S$.
\end{proof}
\end{proposition}

\begin{proposition} \label{prop:intersections-and-chain-unions-of-ideals} \hlt{Intersections and Chain Unions of Ideals}\\
Let $R$ be a ring.
\begin{enumerate}[label=(\roman*)]
\setlength{\itemsep}{0pt}
\item If $\{\mathcal{I}_{\alpha}\}_{\alpha\in A}$ is a nonempty family of ideals of $R$, then $\bigcap_{\alpha\in A}\mathcal{I}_{\alpha}$ is an ideal of $R$.
\item If $\{\mathcal{I}_{\alpha}\}_{\alpha\in A}$ is a chain of ideals of $R$ under inclusion, then $\bigcup_{\alpha\in A}\mathcal{I}_{\alpha}$ is an ideal of $R$.
\end{enumerate}
\begin{proof}
\begin{enumerate}[label=(\roman*)]
\setlength{\itemsep}{0pt}
\item Let $\mathcal{I}=\bigcap_{\alpha\in A}\mathcal{I}_{\alpha}$. If $a,b\in\mathcal{I}$ and $r\in R$, then $a-b,ra,ar\in\mathcal{I}_{\alpha}$ for every $\alpha\in A$, hence $a-b,ra,ar\in\mathcal{I}$.
\item Let $\mathcal{U}=\bigcup_{\alpha\in A}\mathcal{I}_{\alpha}$. If $a,b\in\mathcal{U}$, then $a\in\mathcal{I}_{\alpha}$ and $b\in\mathcal{I}_{\beta}$ for some $\alpha,\beta$. Since the family is a chain, one of $\mathcal{I}_{\alpha},\mathcal{I}_{\beta}$ contains both $a,b$; hence $a-b\in\mathcal{U}$. If $a\in\mathcal{U}$ and $r\in R$, then $a\in\mathcal{I}_{\alpha}$ for some $\alpha$, so $ra,ar\in\mathcal{I}_{\alpha}\subseteq\mathcal{U}$.
\end{enumerate}
\hlt{Structural Commentary}: Chains avoid the usual failure of unions to preserve closure.\\
\hlt{Key Insight}: Comparability places any finite test inside one ideal.
\end{proof}
\end{proposition}
